\documentclass[10pt,a4paper]{amsart}
\usepackage[margin=19mm]{geometry}
\usepackage{amsmath,amssymb,mathtools}
\usepackage[scr=rsfs,cal=euler]{mathalfa}
\usepackage{xfrac}
\usepackage[unicode,hidelinks,bookmarks=false]{hyperref}
\newcommand{\ie}{i.e.\ }
\newcommand{\cf}{cf.\ }
\newcommand{\resp}{respectively\ }
\newcommand{\Subsection}{\subsection}
\providecommand{\nobreakdash}{\nobreak\hbox{-}}
\setlength{\emergencystretch}{2em}
\multlinegap0pt
\usepackage{dsfont}
\usepackage[shortlabels]{enumitem}
\usepackage{amssymb}
\usepackage{bm}
\usepackage{float}
\newtheorem{lemma}{Lemma}
\newtheorem{proposition}{Proposition}
\newcommand{\V}{\mathrm{d} \mathrm{V_g}}
\newcommand{\grad}{\nabla_M}
\newcommand{\red}{\textcolor[rgb]{1.00,0.00,0.00}}
\title{Asymptotics of the principal eigenvalue of an elliptic operator on closed and orientable Riemannian manifolds: small diffusion}
\author{Xin Xu, Kexin Zhang}
\date{Source: arXiv:2608.22181v1 (CC BY 4.0)}
\begin{document}
\maketitle

\noindent\emph{Source and licence.} Asymptotics of the principal eigenvalue of an elliptic operator on closed and orientable Riemannian manifolds: small diffusion. Version arXiv:2608.22181v1 (CC BY 4.0), distributed by arXiv under the Creative Commons Attribution 4.0 International licence, http://creativecommons.org/licenses/by/4.0/, as recorded in the arXiv licence metadata for this exact version and shown on the abstract page https://arxiv.org/abs/2608.22181v1. Full licence notice: \texttt{licenses/arxiv-2608.22181-CC-BY-4.0.txt}.\par\smallskip

\noindent\textit{Editorial scope.} Counted unit: the article's proposition prop:lower (source0.tex lines 1129--1134). The article's complete original proof of it is delivered with it (source0.tex lines 1137--1239, one \texttt{proof} environment of 103 lines). No mathematics is written, repaired or re-derived: the statement and the proof are the article's own lines, in the article's own notation, and no retained line is substituted or re-flowed.  Retained uncounted as the source-local context the proof needs: lines 456--458; lines 483--489; lines 537--542; lines 555--560; lines 1006--1011.  Nothing was omitted from any retained span, so the package carries no omission record.  No retained line contains a citation, so the wrapper carries no bibliography and no bibliography entry.  No retained line contains a figure environment, an image-inclusion command or drawing code, so no image is delivered and the package declares no asset dependency.  Editorial formatting only, in addition: an AMS \texttt{amsart} preamble that restates the article's own declarations and macros used by the retained lines, namely the environments \texttt{lemma}, \texttt{proposition} on separate counters, together with the standard packages those lines need.  The retained environments print in this document as Lemma 1, Proposition 1 in order of appearance, whereas the article prints the same items with the numbering of its own full document.  The complete native source of the article as distributed in the arXiv e-print for this exact version is bundled unchanged as \texttt{sources/208310/source0.tex}.
\begin{equation}\label{Ew}
   -D\Delta_{M}w+\left(\frac{a^2}{4 D}  |\nabla_M f|^2+ \frac{a}{2}\Delta_M f +c\right) w=\lambda(D)w.   
\end{equation}

\begin{lemma}\label{lem:2.1}
For every \(u\in H^1(M,\V)\), we have 
\[
E(u)\ge \frac{a^2}{4D}  \int_{M} \left( | \nabla_M f|^2 +
     \frac{2D}{a} \Delta_M f \right) u^2 \V.
\]
\end{lemma}

\begin{lemma}\label{lem:2.2}
Let \((\lambda(D), w)\) be a solution of \eqref{Ew} and \(\zeta\) be smooth. Define \(W := \zeta w\). Then
\begin{equation}
\int_M \left( D \left| \nabla_M W - \frac{a}{2D} W \nabla_M f \right|^2 + c W^2 \right) \V =\! \lambda(D) \int_M W^2 \V + D \int_M \frac{W^2}{\zeta^2} |\nabla_M \zeta|^2 \V.
\end{equation}
\end{lemma}

\begin{proposition}\label{prop:upper} 
For the principal eigenvalue $\lambda(D)$, we have the following estimate:
\[
\limsup_{D \to 0} \lambda(D) \leq \min_{p \in \Sigma} \left\{ c(p) + \frac{a}{2} \sum_{i=1}^N \left( |\kappa_i(p)| + \kappa_i(p) \right) \right\}.
\]
\end{proposition}

\begin{lemma}\label{lem:2.3}
Let \(p_0 \in M\) and let \(U_R:=B_g(p_0,R)\subset M\) be a geodesic ball contained in a normal neighborhood of \(p_0\), with \(R>0\) sufficiently small. If \(W\in H_0^1(U_R)\) (i.e.,  $W=0$ outside $U_R$), then there exists a constant \(C>0\), depending only on \(f\), \(g\), \(p_0\), and \(N\), such that
\begin{equation}\label{I:llb}
E(W) \geq \frac{a}{2} \left( \sum_{i=1}^N \left( |\kappa_i(p_0)| + \kappa_i(p_0) \right) - C R \right) \int_M W^2 \V.
\end{equation}
\end{lemma}

\begin{proposition}\label{prop:lower} 
For the principal eigenvalue $\lambda(D)$, we have the following estimate:
\[
\liminf_{D \to 0} \lambda(D) \geq \min_{p \in \Sigma} \left\{ c(p) + \frac{a}{2} \sum_{i=1}^N \left( |\kappa_i(p)| + \kappa_i(p) \right) \right\}.
\]
\end{proposition}

\begin{proof}
Recall that the critical set of the Morse function \(f\) on \(M\) is finite. Write
\[
\Sigma=\{p_1,\ldots,p_K\}, \qquad K\in\mathbb N.
\]
Let \(w\) be the principal eigenfunction normalized by \(\int_M w^2 \V=1\). Choose \(R>0\) sufficiently small such that the geodesic balls \(B_g(p_k,R)\) are pairwise disjoint:
\[
B_g(p_k,R)\cap B_g(p_l,R)=\emptyset \qquad (k\neq l).
\]
Set
\[
M_0:=M\setminus \bigcup_{k=1}^K \overline{B_g(p_k,R/2)} .
\]
Then
\[
\{B_g(p_1,R),\ldots,B_g(p_K,R),M_0\}
\]
is an open covering of \(M\).

Let
$\{\zeta_k^2\}_{k=0}^K$
be a smooth partition of unity subordinate to this covering. Thus
\(\zeta_k\in C^\infty(M)\), \(0\leq \zeta_k\leq 1\), and
\[
\sum_{k=0}^K \zeta_k^2(p)=1,
\qquad \forall p\in M,
\]
with
\[
\operatorname{supp} \zeta_0\subset M_0,
\qquad
\operatorname{supp} \zeta_k\subset B_g(p_k,R),
\quad (k=1,\ldots,K).
\]
Moreover, the functions can be chosen so that
\[
|\nabla_M\zeta_k|\leq \frac{C}{R},
\qquad k=0,1,\ldots,K,
\]
for some constant \(C>0\) independent of \(R\).  Since the balls \(B_g(p_k,R)\) are pairwise disjoint, each point \(p\in M\) belongs to at most two sets in the covering; consequently, at most two of the functions \(\zeta_k\) are nonzero at any point. Hence,
\begin{equation}\label{gradient bound}
\sum_{k=0}^K|\nabla_M\zeta_k(p)|^2
\leq\frac{C}{R^2},
\qquad \forall p\in M.
\end{equation}


Define \(W_k = \zeta_k w\) for \(k = 0, 1, \ldots, K\). Then $W_k\in H_0^1(B_g(p_k,R))$ for \(1\leq k\le K\).
Thanks to Lemma \ref{lem:2.3}, we obtain
\[
E(W_k) \geq \frac{a}{2} \left( \sum_{i=1}^N \left( |\kappa_i(p_k)| + \kappa_i(p_k) \right) - C R \right) \int_M W_k^2 \V.
\]
From Lemma \ref{lem:2.2}, for each \(1\le k\le K\),
\begin{align*}
   & \lambda(D) \int_M W_k^2 \V + D \int_M \frac{W_k^2}{\zeta_k^2} |\grad \zeta_k|^2 \V\\=&
E(W_k) +\int_M    c W_k^2   \V \\
\ge&
   \frac{a}{2} \left( \sum_{i=1}^N \left( |\kappa_i(p_k)| + \kappa_i(p_k) \right) - C R \right) \int_M W_k^2 \V
   +\int_M  c W_k^2 \V. 
\end{align*}
Summing over $k=1,...,K$, we get
\begin{align}
&\lambda(D) \sum_{k=1}^K \int_M W_k^2 \V \label{I:lb} \\
\geq& \left( \min_{1 \leq k \leq K} \left\{ c(p_k) + \frac{a}{2} \sum_{i=1}^N (|\kappa_i(p_k)| + \kappa_i(p_k)) \right\}\right)\sum_{k=1}^K \int_M W_k^2 \V  - C R  - D \cdot \frac{C}{R^2}.\nonumber
\end{align}
where the term \(-CR\) comes from Lemma \ref{lem:2.3} and the replacement of \(c(p_k)\) by its local minimum with the estimate \(\sum_{k=1}^K\int_M W_k^2\V\le \int_M w^2\V=1\), while the term \(-C D/R^2\) follows from Lemma \ref{lem:2.2} together with the gradient bound \eqref{gradient bound}
%\red{$\ge c(p)-\epsilon$}

It remains to estimate the contribution from the non-critical region \(M_0\). Since \(M_0\) is compact and contains no critical points of \(f\), there exists \(\delta>0\) such that \(|\grad f| \geq \delta\). Applying Lemma \ref{lem:2.1} to \(W_0\), we obtain
\begin{align*}
E(W_0)\ge &\frac{a^2}{4D}  \int_{M} \left[ | \nabla_M f|^2 +
     \frac{2D}{a} \Delta_M f \right] W_0^2 \mathrm{d} \mathrm{V_g}\\
     =&\frac{a^2}{4D}  \int_{M_0} \left[ | \nabla_M f|^2 +
     \frac{2D}{a} \Delta_M f \right] W_0^2 \mathrm{d} \mathrm{V_g}\\
     \ge& \frac{a^2}{4D}  \int_{M_0} \left[ \delta^2 -
     \frac{2D}{a} \|\Delta_M f \|_{L^\infty(M)}\right] W_0^2 \mathrm{d} \mathrm{V_g}.
\end{align*}
For $D>0$ sufficiently small such that 
%$0<D<\frac{a\delta^2}{4\|\Delta_M f \|_{L^\infty(M)}}$, 
\(\frac{2D}{a}\|\Delta_M f\|_{L^\infty(M)} \le \delta^2/2\), we have
\begin{align*}
\int_{\Omega_0}  W_0^2 \mathrm{d} \mathrm{V_g}\leq
\frac{4D}{a^2\left[ \delta^2 -
     \frac{2D}{a} |\Delta_M f |_{L^\infty}\right]}E(W_0)\leq \frac{8 D}{\delta^2 a^2}E(W_0).
\end{align*}
On the other hand, from Lemma \ref{lem:2.2} and the upper bound for $\lambda(D)$ above (Proposition \ref{prop:upper}), we have 
\begin{align*}
E(W_0)\leq |\lambda(D)|\int_M w^2\V+\max_{p\in M}c(p)\int_M w^2\V+C_R D\int_M w^2\V \leq C_R
\end{align*}
for some constant \(C_R>0\) depending on \(R\) (and on \(f,g\)) but independent of \(D\). Hence,    
\[
\int_Mw^2 \sum_{i=1}^K \zeta_k^2\V=1-\int_M W_0^2 \V\ge 1-C_R D.
\]
Letting \(D\to0\) (with \(R>0\) fixed), we obtain
\[
\int_Mw^2 \sum_{i=1}^K \zeta_k^2\V=\sum_{k=1}^K \int_M W_k^2 \V \to 1.
\]
Taking \(\liminf\) as \(D \to 0\) in \eqref{I:lb} and then letting \(R \to 0\), we conclude
\[
\liminf_{D \to 0} \lambda(D) \geq \min_{p \in \Sigma} \left\{ c(p) + \frac{a}{2} \sum_{i=1}^N (|\kappa_i(p)| + \kappa_i(p)) \right\},
\]
which completes the proof.  
\end{proof}

\end{document}
